% Part: applied-modal-logics
% Chapter: epistemic-logic
% Section: introduction

\documentclass[../../../include/open-logic-section]{subfiles}

\begin{document}

\olfileid{aml}{el}{int}

\olsection{جان پہچان}

جیویں امکانی منطق \emph{امکانی قضیاں} تے اوہناں دے آپسی استلزام
دے تعلقاں نوں ویکھدی اے، اویویں علم دی منطق \emph{علم بارے قضیاں}
تے اوہناں دے آپسی استلزام دے تعلقاں نوں ویکھدی اے۔ ایتھے اک رکنی
رابطیاں دی تعبیر علم یا منّے ہوئے خیال دے لحاظ نال ہوندی اے، تاں جو
جانن تے منّن دے نمونے بن سکن۔ عام امکانی تعبیر وچ اوہی عامل امکان
تے لزوم دسدے نیں۔ مثلاً اسیں ایہو جیہے بیان آکھنا چاہ سکدے آں:

\begin{enumerate}
\item رچرڈ نوں پتا اے کہ کیلگری البرٹا وچ اے۔
\item آڈری سمجھدی اے کہ ہو سکدا اے کوئی کُتا صوفے اُتے ہووے۔
\item رچرڈ نوں پتا اے کہ آڈری نوں پتا اے کہ اوہدی جماعت منگل والے دن ہوندی اے۔
\item ہر کسے نوں پتا اے کہ اک سال وچ 12 مہینے ہوندے نیں۔
\end{enumerate}
اجوکی علم دی منطق دا سرا اکثر یاکو ہنٹکّا دی 1962 دی کتاب
\emph{Knowledge and Belief} نال جوڑیا جاندا اے۔ ایہ اوس ویلے لکھی گئی
جدوں منطق وچ ممکن جہاناں دیاں معنیات دی ورتوں ودھ رہی سی۔ علم دیاں
منطقاں ہور امکانی منطقاں والے زیادہ تر معنیاتی وسیلے ای ورتدیاں نیں،
پر اوہناں دی تعبیر وکھری اے۔ وڈی تبدیلی ایہ اے کہ اسیں \emph{رسائی دے تعلق}
نوں کیہدا نشان سمجھدے آں۔ علم دیاں منطقاں وچ ایہ تعلق \emph{علم دے لحاظ
نال امکان} دی کوئی صورت وکھاؤندے نیں۔ اسیں ویکھاں گے کہ علم دا جیہڑا
خیال اسیں نمونے وچ لیندے آں، اوہ رسائی دے تعلق اُتے لاؤن والیاں قیداں
نوں کِویں بدلدا اے۔ اسیں ایہ وی ویکھاں گے کہ علم دے لحاظ نال تعبیر
کرن تے فریم خاصیتاں تے فارمولیاں دے میل دا نظریہ کیویں ہوندا اے۔
اُتلیاں مثالاں وچ دو فاعل، رچرڈ تے آڈری، تے اوہناں دونہاں دی جانکاری
دا آپسی تعلق اے۔ اسیں کئی فاعلاں والیاں علم دیاں منطقاں ویکھاں گے،
جیہناں وچ ایہو جیہیاں گلاں آکھیاں جا سکدیاں نیں۔ اک فاعل والی علم دی
منطق صرف اک فرد دے جانن یا منّن بارے گل کردی اے۔

\end{document}
